\section{Implementation of LangGraph}

\subsection{Web Search Branch}

The web research branch represents a critical component of the implemented information retrieval system, designed to gather, synthesize, and validate web-based knowledge at scale. This branch operates within the LangGraph-based agentic architecture, implementing an iterative refinement cycle that combines LLMs, vector-based semantic search, and quality evaluation metrics to produce coherent, evidence-grounded research summaries.

The web research branch is implemented as a specialized workflow within a larger knowledge production pipeline. The system operates through a sequence of interconnected nodes, each performing distinct transformations on a shared state object (\texttt{SummaryState}). The architecture leverages asynchronous execution patterns to maximize throughput and responsiveness across the retrieval-generation-evaluation cycle.

Key architectural principles include state-driven execution, where all operations maintain and transform a centralized state object containing research topics, accumulated queries, search results, and evolving summaries; asynchronous pipeline design, where each node executes as an async coroutine enabling non-blocking operations and integration with long-running external services; tool composition, allowing multiple search backends (Tavily, Google Search, DuckDuckGo) to be composed through a unified agent interface; and checkpoint persistence for recovery and state inspection.

The implementation relies on several key frameworks and libraries: \textbf{LangGraph} for state machine and workflow orchestration, \textbf{LangChain} for document loading and embedding management, \textbf{Ollama} for local language model serving, \textbf{Chroma} for vector database functionality, \textbf{BeautifulSoup} for HTML parsing, \textbf{DeepEval} for quality assessment metrics, and \textbf{CopilotKit} for real-time UI integration.

\subsubsection{Query Generation}

The workflow begins with a user-provided research topic. Rather than immediately executing a query, the system first generates a refined, structured search query through an LLM-guided process:

\begin{lstlisting}
1. Input: research\_topic (string)
2. Format research topic with prompt template
3. Invoke ChatOllama with JSON mode (temperature=0)
4. Parse returned JSON to extract 'query' field
5. Output: search\_query (string)
\end{lstlisting}

This generation step provides three key benefits. First, setting temperature to 0 ensures deterministic query formulation, guaranteeing that the same research topic produces nearly identical search queries for reproducibility. Second, enforcing JSON mode constrains the LLM output to structured, parseable results rather than free-form text, preventing malformed queries. Third, the LLM leverages its domain knowledge to reformulate natural language research topics into effective search terminology that maximizes retrieval relevance.

\subsubsection{Web Search Execution}

Once a query is formulated, the system executes a web search using a configurable agent architecture. Search parameters include maximum results of "x" documents per search with full document content retained, using the Tavily search API with fallback options to Google/DuckDuckGo:

\begin{lstlisting}
1. Create Agent with Ollama backend (qwen3:latest)
2. Attach TavilyTools for web search capability
3. Construct search instruction prompt + query
4. Execute agent.arun() in non-blocking mode
5. Parse response content as search results
6. Emit progress to UI
\end{lstlisting}

The use of \texttt{asyncio.to\_thread()} is critical as it prevents blocking the event loop during the agent's synchronous execution, maintaining responsiveness in the overall workflow. Search results are collected as unstructured text containing webpage snippets, titles, and URLs, stored in state as \texttt{raw\_search\_result} for subsequent parsing.

\subsubsection{Result Parsing and Formatting}

After search execution, raw results are parsed into a structured format using a second LLM pass:

\begin{lstlisting}
1. Input: raw\_search\_result (unstructured text)
2. Invoke ChatOllama with web\_search\_expected\_output prompt
3. Extract JSON from response
4. Validate and normalize source information
\end{lstlisting}

After JSON extraction, \textbf{normalization} standardizes source metadata (title, URL, date) into a consistent schema. This processing stage reduces context bloat while preserving informational diversity, balancing conciseness with coverage quality for downstream summarization.



\subsubsection{Iterative Summarization}

For each batch of search results, the system either generates a new summary or extends an existing one. This two-stage approach allows the system to accumulate knowledge progressively across multiple search iterations. Output is post-processed to remove any XML-like thinking tags from the language model, ensuring a clean, presentation-ready summary.

The system uses a two-stage summarization approach:

\begin{lstlisting}
if existing_summary exists:
Extend existing summary with new search results
else:
Generate new summary from search results

Invoke ChatOllama and remove intermediate reasoning tags
\end{lstlisting}

This enables progressive knowledge accumulation across multiple search iterations, with post-processing to clean any intermediate reasoning from the model output.

\subsubsection{Reflection and Follow-up Query Generation}

After each summary generation, the system reflects on gaps in coverage to guide subsequent searches:

\begin{lstlisting}
1. Invoke ChatOllama with reflection_instructions prompt
2. Input: current running\_summary + research\_topic
3. Use JSON mode to extract structured follow-up query
4. Validate that follow-up\_query is not null
5. Fallback to generic query if generation fails
\end{lstlisting}

The reflection phase ensures that the system doesn't re-query the same information but instead targets unexplored dimensions of the research topic, enabling progressive exploration and preventing query stagnation. The LLM identifies knowledge gaps by analyzing the current summary against the original research topic, with the resulting follow-up\_query used in the next iteration of web search.

\subsubsection{Loop Control and Termination}

The system maintains a \texttt{research\_loop\_count} to prevent infinite loops:

\begin{lstlisting}
if research\_loop\_count <= max\_web\_research\_loops:
    route = "web_research"  # Continue another iteration
else:
    route = "finalize_summary"  # Proceed to evaluation
\end{lstlisting}

The \texttt{max\_web\_research\_loops} parameter is configurable per execution, allowing flexible control over research depth.

\subsubsection{Quality Evaluation}

Upon reaching the loop limit, the system enters an evaluation phase using four complementary RAG quality metrics:

\begin{enumerate}
    \item \textbf{Faithfulness} : Measures whether summary claims are grounded in retrieved documents, preventing hallucinated information not present in sources.
    
    \item \textbf{Answer Relevancy} : Assesses whether summary content directly addresses the research topic, ensuring topical alignment with the original query.
    
    \item \textbf{Contextual Relevancy}: Evaluates whether retrieved documents are genuinely relevant to the context, reducing noise from tangentially related search results.
    
    \item \textbf{Hallucination Detection} : Identifies unsupported claims in the summary, flagging content fabrication risk.
\end{enumerate}


\subsubsection{Summary Finalization}

Upon completing iterations and evaluation, the system generates a final deliverable:

\begin{lstlisting}
1. Collect all accumulated sources: sources_gathered
2. Merge into structured format:
   - Summary section: running_summary
   - Sources section: formatted_sources list
3. Return with evaluation_results metadata
4. Emit exit signal to CopilotKit UI
\end{lstlisting}

All sources are tracked throughout execution and formatted consistently, preserving document title and URL, extraction timestamp, relevance confidence scores, and token budget information. This ensures reproducibility and enables users to trace claims back to original sources.

\subsubsection{Error Handling and Robustness}

The system implements multi-stage JSON parsing with fallback strategies:

\begin{lstlisting}
1. Attempt direct json.loads(content)
2. If fails: extract JSON substring using regex
3. If fails: attempt markdown fence removal then retry
4. Fallback: return empty/default result with error logging
\end{lstlisting}

This defensive approach handles variations in LLM output formats across different model variants and API versions.

All async functions include try-except blocks that log errors with context information, emit error status to the UI, return safe default states to prevent pipeline stalls, and preserve partial results when possible.

GPU/CUDA resources are explicitly managed through \texttt{torch.cuda.empty\_cache()} calls after embedding operations, proper vectorstore connection closure, and memory release between batch operations.

\subsubsection{State Management and Performance}

The central state object (\texttt{SummaryState}) maintains research topic, search query, raw search results, formatted source metadata, accumulated raw results, iteratively refined summary, iteration counter, and evaluation results. State transitions follow strict ordering enforced by the graph topology, ensuring that operations occur in valid sequences.

The web research branch implements a principled, iterative approach to knowledge synthesis combining keyword search, semantic retrieval, language model reasoning, and formal quality assessment. By maintaining clear separation of concerns across query generation, search, parsing, summarization, reflection, and evaluation phases, the system achieves both robustness and flexibility for diverse research applications.

\subsection{Academic Research Branch}

The academic research branch automates rigorous retrieval, synthesis, and evaluation of scientific literature by integrating scholarly and preprint search tools within the agentic LangGraph workflow. This branch complements web-based research by ensuring coverage of peer-reviewed and domain-specific academic sources, thus improving both the reliability and depth of generated knowledge.

\subsubsection{Academic Query Generation}

The workflow initiates with LLM-guided query construction, where a natural language research topic is transformed into a structured, field-focused query suitable for academic databases:

\begin{lstlisting}

Input: research_topic (string)

Format research topic with prompt template 

Invoke ChatOllama with JSON mode (temperature=0)

Parse returned JSON to extract 'query' field

Emit status to UI via CopilotKit

Output: search_query (string)
\end{lstlisting}

This ensures deterministic, reproducible query formulation by constraining the LLM to JSON-structured output. The system reformulates natural language research topics into optimized academic search terminology while guaranteeing consistent, parseable results across multiple executions.

\subsubsection{Academic Literature Retrieval}

Once a query is formulated, the system executes targeted searches across multiple academic databases using the agentic architecture:

\begin{lstlisting}

Create Agent with Ollama backend

Attach ArxivTools for academic preprint search

Construct academic search instruction prompt + query

Execute agent.arun() in non-blocking mode

Parse response content as academic results

Emit progress to UI

Output: raw_academic_result (unstructured academic metadata)
\end{lstlisting}

The system targets multiple academic repositories, particularly Arxiv for preprints and via agent tool composition, Semantic Scholar and Google Scholar APIs, aggregating results including titles, abstracts, author lists, publication venues, and metadata. Results are collected as unstructured text containing paper summaries, author information, and citations, stored in state for subsequent parsing.

\subsubsection{Source Parsing and Deduplication}

After search execution, raw results are parsed into a structured format using a second LLM pass:

\begin{lstlisting}

Input: raw_academic_result (unstructured academic output)

Invoke ChatOllama with academic_expected_output prompt

Extract JSON from response (with robustness for variations)

Validate and normalize academic metadata

\end{lstlisting}



\subsubsection{Summarization of Academic Findings}

For each batch of academic sources, the system generates or extends a running summary:

\begin{lstlisting}
if existing_summary exists:
human_message = "Extend the existing summary: {existing_summary}
Include new academic sources: {academic_sources}
That addresses the research topic: {research_topic}"
else:
human_message = "Generate a summary of these academic sources:
{academic_sources}
That addresses the research topic: {research_topic}"

invoke ChatOllama with:
- system_prompt: academic_summarizer_instructions
- user_message: human_message

remove_think_tags(response) # Clean any intermediate reasoning
\end{lstlisting}

This two-stage approach enables progressive knowledge accumulation across multiple search iterations. Post-processing removes any XML-like thinking tags inserted by the language model, ensuring clean, presentation-ready academic summaries.

\subsubsection{Iterative Reflection and Query Refinement}

After each summary, the system reflects on gaps in coverage to guide subsequent searches:

\begin{lstlisting}

Invoke ChatOllama with reflection_instructions prompt

Input: current running_summary + research_topic

Use JSON mode to extract structured follow-up query

Validate that follow-up_query is not null

Fallback to generic academic query if generation fails
\end{lstlisting}

The reflection phase ensures that the system doesn't re-query the same papers but instead targets unexplored research dimensions or methodological gaps, enabling progressive exploration and preventing query stagnation.

\subsubsection{Loop Control and Branch Termination}

The system maintains a \texttt{research\char`_loop\char`_count}
 strictly limiting the number of academic exploration cycles:

\begin{lstlisting}
if research_loop_count <= max_academic_research_loops:
route = "academic_research" # Continue another iteration
else:
route = "finalize_academic_summary" # Proceed to evaluation
\end{lstlisting}

The conservative loop limit balances coverage depth with computational efficiency, as academic literature searches typically encounter diminishing returns more rapidly than web searches.

\subsubsection{Quality Evaluation of Academic Summaries}

Upon reaching the loop limit, the system enters an evaluation phase using the same four complementary RAG quality metrics:

\begin{enumerate}
\item \textbf{Faithfulness} : Ensures summary claims are grounded in retrieved academic papers, preventing hallucinated citations.

text
\item \textbf{Answer Relevancy} : Verifies that academic synthesis directly addresses the research topic.

\item \textbf{Contextual Relevancy} : Evaluates whether retrieved papers are genuinely topically relevant.

\item \textbf{Hallucination Detection} : Identifies unsupported claims or fabricated references.
\end{enumerate}

The evaluation process executes as follows:

\begin{lstlisting}

User prompt: "Do you want to evaluate this academic summary?"

If yes:
a. Run evaluation in separate thread via asyncio.to_thread()
b. Create LLMTestCase with:
- input: research_topic
- actual_output: running_summary
- context: academic_sources (full list)
- retrieval_context: academic_sources (for Faithfulness)
c. Initialize Ollama model (deepseek-r1:latest)
d. Invoke deepeval.evaluate() with all four metrics
e. Extract metric scores (range 0.0--1.0)
f. Emit results to UI with per-metric scores

If no:
a. Skip evaluation
b. Proceed to finalization
\end{lstlisting}

The evaluation runs in a separate thread to avoid blocking uvloop (the async event loop), critical because DeepEval's metrics perform LLM calls internally that would otherwise stall the entire pipeline.

\subsubsection{Summary Finalization and Output Structuring}

Upon completing iterations and evaluation, the system generates a final deliverable:

\begin{lstlisting}

Collect academic sources: academic_sources_gathered

Merge into structured format:

Academic Summary section: running_summary

References section: formatted_sources list with metadata

Evaluation Results section: metric scores (if evaluated)

Return with evaluation_results metadata

Emit exit signal to CopilotKit UI
\end{lstlisting}

All sources are tracked throughout execution and formatted consistently, preserving title, authors, venue, DOI/URL, publication date, and token budget information. This ensures reproducibility and enables users to trace claims back to original academic sources with complete bibliographic information.


\subsection{Code Generation Branch}

The Code Generation Branch is a specialized subsystem designed to orchestrate the generation, validation, and iterative refinement of executable Python code in response to research queries. It combines RAG, multi-pass search optimization, and sandboxed code validation to produce semantically coherent, production-ready code artifacts. The branch operates within the LangGraph's architecture and integrates specialized language models, vector databases, and secure execution environments to ensure code quality and correctness.

\subsubsection{Architecture and Multi-Agent Code Generation}

The Code Generation Branch follows a modular, multi-agent architecture with distinct components:

\begin{itemize}
  \item \textbf{Orchestrator Agent}: Routes user queries to specialized agents based on semantic analysis
  \item \textbf{Specialized Code Agents}: Domain-specific agents (snnTorch, NNI) that generate code with repository-specific context
  \item \textbf{Vector Retrieval System}: Two dedicated vectorstores with CodeBERT embeddings
  \item \textbf{Code Parser and Extractor}: Handles multi-file code generation using explicit file markers
  \item \textbf{Sandbox Executor}: E2B-based secure execution environment for validation
  \item \textbf{Feedback Processing}: User-guided refinement loop with decision classification
\end{itemize}

The system enforces strict separation between code generation (via LLM) and code validation (via sandboxed execution), enabling error detection before delivery. Generated code uses explicit file markers (\texttt{\# FILE: filename.py}) for modular composition, with each file representing a logical component.

\paragraph{Orchestrator Agent Workflow}
Algorithm~\ref{alg:orchestrator} describes the orchestrator's decision logic for routing queries to specialized agents.

\begin{algorithm}
\caption{Orchestrator Agent Tool Selection}
\label{alg:orchestrator}
\begin{algorithmic}[1]
\REQUIRE User query \(q\), available tools \(T = \{\text{snnTorch}, \text{NNI}\}\)
\ENSURE Selected tools with focused queries
\STATE \(prompt \gets \) Create tool selection prompt with query \(q\)
\STATE \(response \gets \text{LLM}(prompt)\)
\STATE \(toolCalls \gets \text{JSON.parse}(response)\)
\FOR{each \(call \in toolCalls\)}
    \STATE \(toolName \gets call.\text{name}\)
    \STATE \(toolQuery \gets call.\text{query}\)
    \STATE \(result \gets \text{ExecuteTool}(toolName, toolQuery)\)
    \IF{\(result.\text{error} = \text{null}\) \AND \(|result.\text{code}| > 50\)}
        \STATE \(validOutputs.\text{add}(toolName, result)\)
    \ENDIF
\ENDFOR
\RETURN \(validOutputs\)
\end{algorithmic}
\end{algorithm}

The orchestrator validates each tool output before acceptance by checking: (1) the error field is null, (2) code length exceeds 50 characters, and (3) file markers are present.

\subsubsection{Vectorstore Initialization and Query Optimization}

\paragraph{Dual Vectorstore Configuration}
The system maintains two specialized vectorstores implemented using Chromadb with persistent storage:

\begin{itemize}
  \item \textbf{snnTorch Vectorstore}: Online documentation from \texttt{snntorch.readthedocs.io}, recursive loading with custom depth limit, focuses on spiking neural networks, LIF neurons, training methods
  \item \textbf{NNI Vectorstore}: Local examples from \texttt{esempi\_NNI/} directory plus API reference document, focuses on hyperparameter tuning configurations and experiment setup patterns
\end{itemize}

Both employ CodeBERT embeddings (\texttt{mchochlov/codebert-base-cd-ft}), optimized for code semantics with cosine similarity metric. The vectorstore initialization follows an equal pattern (Algorithm~\ref{alg:vectorstore-init}).

\begin{algorithm}
\caption{Vectorstore Initialization}
\label{alg:vectorstore-init}
\begin{algorithmic}[1]
\REQUIRE Persist directories \(D_{snn}, D_{nni}\), source URLs/paths
\ENSURE Initialized vectorstores
\IF{\(D_{snn}\) exists \AND \(D_{nni}\) exists}
    \STATE Load existing vectorstores from disk
    \RETURN cached vectorstores
\ENDIF
\STATE \(docs_{snn} \gets \text{RecursiveURLLoader}(\text{snnTorch docs}, \text{depth}=7)\)
\STATE \(docs_{nni} \gets \text{DirectoryLoader}(\text{esempi\_NNI}) + \text{APIRef}\)
\STATE \(docs_{nni}.\text{insert}(0, \text{APIRef})\) \COMMENT{Priority insertion}
\STATE \(chunks_{snn} \gets \text{TextSplitter}(docs_{snn}, \text{chunk}=512, \text{overlap}=50)\)
\STATE \(chunks_{nni} \gets \text{TextSplitter}(docs_{nni}, \text{chunk}=512, \text{overlap}=50)\)
\STATE \(V_{snn} \gets \text{Chroma.from\_documents}(chunks_{snn}, \text{CodeBERT})\)
\STATE \(V_{nni} \gets \text{Chroma.from\_documents}(chunks_{nni}, \text{CodeBERT})\)
\STATE Persist \(V_{snn}\) to \(D_{snn}\), \(V_{nni}\) to \(D_{nni}\)
\RETURN \(V_{snn}, V_{nni}\)
\end{algorithmic}
\end{algorithm}

\paragraph{CodeBERT Memory Optimization}
To handle large-scale embedding on constrained GPU memory (8GB VRAM), a custom batching strategy is implemented:

\begin{verbatim}
batch_size = 8
for i in range(0, len(texts), batch_size):
    batch = texts[i:i+batch_size]
    embeddings = model.encode(batch)
    gc.collect()
    torch.cuda.empty_cache()
\end{verbatim}

This approach trades computational speed (15-20\% overhead) for memory safety, enabling vectorstore creation on consumer hardware.

\paragraph{LLM-Guided Query Generation and Ranking}
Rather than static keyword queries, the system dynamically generates a custom number of diverse search queries tailored to each research question using a two-phase LLM process :

\textbf{Phase 1: Query Generation}
\begin{verbatim}
Generate 6 DIVERSE, SPECIFIC search queries covering:
- Architecture and model structure
- Parameters and configuration
- Training methodology
- Implementation patterns
- Evaluation metrics
Return JSON: ["query1", "query2", ..., "query6"]
\end{verbatim}

\textbf{Phase 2: Query Ranking}
\begin{verbatim}
User Question: {question}
Candidate Queries: {generated_queries}

Rank by RELEVANCE to answering the question.
Return JSON: {
  "ranked_queries": ["most_relevant", "second", ...],
  "reasoning": "explanation"
}
\end{verbatim}

Ranked queries receive relevance weights: \(w_i = 1.0 - (i-1) \times 0.05\) for query rank \(i\). Retrieved documents are tagged with their source query and weight, prioritizing high-relevance sources during LLM context construction.

\paragraph{Multi-Pass Retrieval with Weighted Context}
Algorithm~\ref{alg:multipass-retrieval} implements the weighted multi-pass retrieval strategy.

\begin{algorithm}
\caption{Multi-Pass Weighted Retrieval}
\label{alg:multipass-retrieval}
\begin{algorithmic}[1]
\REQUIRE Ranked queries \(Q = [q_1, q_2, \ldots, q_n]\), vectorstore \(V\), \(k=20\)
\ENSURE Weighted context string
\STATE \(context \gets \text{""}\)
\STATE \(totalDocs \gets 0\)
\FOR{\(i \gets 1\) to \(n\)}
    \STATE \(w_i \gets 1.0 - (i-1) \times 0.05\)
    \STATE \(docs \gets V.\text{retrieve}(q_i, k)\)
    \STATE \(totalDocs \gets totalDocs + |docs|\)
    \FOR{each \(doc \in docs\)}
        \STATE \(source \gets doc.\text{metadata.source}\)
        \STATE \(content \gets doc.\text{page\_content}[0:600]\)
        \STATE \(context \gets context + \text{f"Query \{i\} - \{source\} - weight \{w\_i:.2f\}: \{content\}---"}\)
    \ENDFOR
\ENDFOR
\RETURN \(context[0:8000]\) \COMMENT{Limit to 8k chars}
\end{algorithmic}
\end{algorithm}

The weighted context enables the LLM to prioritize information from higher-ranked queries during code generation.

\subsubsection{Specialized Agents and Unified Response Parsing}

\paragraph{snnTorch Agent}
The snnTorch agent specializes in generating spiking neural network code. Its output schema guarantees:

\begin{verbatim}
{
  "code": "# FILE: model.py\n...\n# FILE: utils.py\n...",
  "files": {"model.py": "...", "utils.py": "..."},
  "summary": "Generated SNN model with LIF neurons",
  "confidence": 0.95,
  "queries_used": 6,
  "error": null
}
\end{verbatim}

The agent constructs prompts emphasizing documentation priority: (1) retrieved context (primary), (2) snnTorch official docs (gap-filling), (3) general PyTorch knowledge (fallback).

\paragraph{NNI Agent}
The NNI agent generates hyperparameter tuning configurations with strict JSON output requirements. It enforces dual-file patterns (\texttt{config.py} for search space, \texttt{train.py} for training loop) and validates JSON parseability before returning results.

\paragraph{Unified Response Parser}
Algorithm~\ref{alg:unified-parser} implements a multi-strategy parsing approach with graceful degradation.

\begin{algorithm}
\caption{Unified Response Parser}
\label{alg:unified-parser}
\begin{algorithmic}[1]
\REQUIRE LLM response text \(R\)
\ENSURE Parsed result dict
\STATE \(R \gets R.\text{strip}()\)
\IF{\(R\) contains markdown code fences}
    \STATE Remove fences: \(R \gets R.\text{remove\_fences}()\)
\ENDIF
\STATE \textbf{Strategy 1: JSON Parsing}
\STATE \textbf{try}
    \STATE \(data \gets \text{JSON.parse}(R)\)
    \STATE Map alternative field names to canonical \texttt{code} field
    \STATE Normalize escape sequences: \(data[\text{"code"}] \gets \text{normalize}(data[\text{"code"}])\)
    \RETURN \(data\)
\STATE \textbf{catch} JSONDecodeError
    \STATE Continue to Strategy 2
\STATE \textbf{end try}
\STATE \textbf{Strategy 2: Markdown Code Block Extraction}
\STATE \(blocks \gets \text{REGEX.findall}(\text{r'```[a-z]*\textbackslash n(.*?)```'}, R)\)
\IF{\(|blocks| > 0\)}
    \STATE Deduplicate and merge blocks
    \RETURN \(\{\text{"code"}: \text{merged}, \text{"parse\_method"}: \text{"markdown"}\}\)
\ENDIF
\STATE \textbf{Strategy 3: Raw Text Fallback}
\IF{\(|R| > 50\)}
    \RETURN \(\{\text{"code"}: R, \text{"parse\_method"}: \text{"raw"}\}\)
\ENDIF
\RETURN \(\{\text{"error"}: \text{"Could not parse response"}\}\)
\end{algorithmic}
\end{algorithm}

\paragraph{Escape Sequence Normalization}
LLM responses often contain literal escape sequences from JSON or markdown encoding. The normalizer handles:

\begin{verbatim}
def normalize_code_escapes(text):
    text = text.replace(r'\n', '\n')   # Literal \n → newline
    text = text.replace(r'\t', '\t')   # Literal \t → tab
    text = text.replace(r'\"', '"')    # Escaped quotes
    text = text.replace(r'\\', '\\')   # Double backslash
    return text
\end{verbatim}

This ensures code readability and correct execution in the sandbox.

\paragraph{File Extraction from Code}
Algorithm~\ref{alg:file-extraction} extracts individual files from the unified code string.

\begin{algorithm}
\caption{File Extraction with Regex Markers}
\label{alg:file-extraction}
\begin{algorithmic}[1]
\REQUIRE Code string \(C\) with \texttt{\# FILE:} markers
\ENSURE Dict mapping filename → content
\STATE \(pattern \gets \text{r'\textasciicircum\# FILE:\textbackslash s*(\textbackslash S+\textbackslash .py)\textbackslash s*\$'}\)
\STATE \(lines \gets C.\text{split}('\textbackslash n')\)
\STATE \(files \gets \{\}\), \(currentFile \gets \text{null}\), \(currentContent \gets []\)
\FOR{each \(line \in lines\)}
    \STATE \(match \gets \text{REGEX.match}(pattern, line)\)
    \IF{\(match \neq \text{null}\)}
        \IF{\(currentFile \neq \text{null}\)}
            \STATE \(files[currentFile] \gets \text{join}(currentContent)\)
        \ENDIF
        \STATE \(currentFile \gets match.\text{group}(1)\)
        \STATE \(currentContent \gets []\)
    \ELSIF{\(currentFile \neq \text{null}\)}
        \STATE \(currentContent.\text{append}(line)\)
    \ENDIF
\ENDFOR
\IF{\(currentFile \neq \text{null}\)}
    \STATE \(files[currentFile] \gets \text{join}(currentContent)\)
\ENDIF
\RETURN \(files\) if \(files \neq \{\}\) else \(\{\text{"main.py"}: C\}\)
\end{algorithmic}
\end{algorithm}

\paragraph{Multi-Agent Code Integration}
When multiple agents contribute code (e.g., snnTorch + NNI), intelligent merging applies a priority scheme:

\begin{enumerate}
  \item \textbf{Priority 1}: \texttt{config.py} (configuration)
  \item \textbf{Priority 2}: \texttt{search\_space.json} (NNI search space)
  \item \textbf{Priority 3}: \texttt{model.py} (architecture)
  \item \textbf{Priority 4}: \texttt{utils.py} (helpers)
  \item \textbf{Priority 5}: \texttt{train.py} (training loop)
  \item \textbf{Priority 6}: \texttt{main.py} (entry point)
\end{enumerate}

Files are sorted by priority and concatenated with file markers. Duplicates are resolved by keeping the first occurrence (higher-priority agent).

\subsubsection{Sandboxed Validation and Feedback Loop}

\paragraph{E2B Sandbox Execution}
Generated code is executed in an isolated E2B sandbox (\texttt{code-interpreter-v1} template) with 1GB memory and CPU-only execution. Algorithm~\ref{alg:sandbox-exec} describes the multi-file execution process.

\begin{algorithm}
\caption{Sandboxed Multi-File Execution}
\label{alg:sandbox-exec}
\begin{algorithmic}[1]
\REQUIRE Code string \(C\) with file markers
\ENSURE Execution results
\STATE \(files \gets \text{parse\_multifile\_code}(C)\)
\STATE \(mainFile \gets \text{detect\_main\_file}(files)\) \COMMENT{Priority: main.py > train.py > experiment.py}
\STATE \(sandbox \gets \text{E2B.create}(\text{"code-interpreter-v1"})\)
\STATE \(packages \gets \text{extract\_packages}(C)\) \COMMENT{Parse import statements}
\STATE \(sandbox.\text{run}(\text{"pip install "} + \text{join}(packages))\)
\FOR{each \((filename, content) \in files\)}
    \STATE \(sandbox.\text{upload}(\text{"/home/user/"} + filename, content)\)
\ENDFOR
\STATE \(result \gets sandbox.\text{run}(\text{"python "} + mainFile)\)
\STATE \(sandbox.\text{kill}()\)
\RETURN \(\{
    \text{"exit\_code"}: result.\text{exit\_code},
    \text{"stdout"}: result.\text{stdout},
    \text{"stderr"}: result.\text{stderr}
\}\)
\end{algorithmic}
\end{algorithm}

\paragraph{Dependency Auto-Detection and Installation}
Import statements are extracted using regex:

\begin{verbatim}
pattern = r"(?:^import|^from)\s+(\w+)"
matches = re.findall(pattern, code, re.MULTILINE)
\end{verbatim}

A mapping table handles non-standard package names:

\begin{verbatim}
pip_mapping = {
    'cv2': 'opencv-python',
    'PIL': 'Pillow',
    'sklearn': 'scikit-learn',
    'snntorch': 'snntorch',
}
\end{verbatim}

Built-in modules (\texttt{os}, \texttt{sys}, \texttt{json}, etc.) are filtered out before installation.

\paragraph{Static Type Analysis}
Before execution, Pyright analysis detects type errors, undefined variables, and incompatible method calls. Results are non-blocking but reported to users for transparency.

\paragraph{Interactive Feedback Loop}
After code generation or sandbox execution, users select from four actions via an interrupt mechanism:

\begin{itemize}
  \item \textbf{approve}: Accept code and terminate
  \item \textbf{regenerate}: Request modifications with feedback
  \item \textbf{evaluate}: Compare against reference implementation
  \item \textbf{execute}: Run specific file in sandbox
\end{itemize}

Algorithm~\ref{alg:feedback-classification} describes the LLM-based feedback classification.

\begin{algorithm}
\caption{Feedback Decision Classification}
\label{alg:feedback-classification}
\begin{algorithmic}[1]
\REQUIRE User feedback text \(F\)
\ENSURE Action and optional filename
\STATE \(prompt \gets \text{Create classification prompt with examples}\)
\STATE \(response \gets \text{LLM}(prompt + F)\)
\STATE \(data \gets \text{JSON.parse}(response)\)
\STATE \(action \gets data[\text{"response"}]\) \COMMENT{approve/regenerate/evaluate/execute}
\STATE \(filename \gets data.\text{get}(\text{"file\_name"}, \text{null})\)
\RETURN \((action, filename)\)
\end{algorithmic}
\end{algorithm}

Example feedback classifications:

\begin{verbatim}
"execute train.py" → {"response": "execute", "file_name": "train.py"}
"Make it faster" → {"response": "regenerate"}
"Looks good" → {"response": "approve"}
"Compare with reference" → {"response": "evaluate"}
\end{verbatim}

\paragraph{Evaluation Mode and Performance Metrics}
When users select \textit{evaluate}, the system enters comparison mode with code normalization and performance measurement. Normalization standardizes variable names, imports, and formatting while preserving logic. Performance metrics are injected via wrapper code:

\begin{verbatim}
import time
import torch

start = time.time()
output = model(input_tensor)
elapsed = time.time() - start

params = sum(p.numel() for p in model.parameters() if p.requires_grad)
peak_mem = torch.cuda.max_memory_allocated() / 1024 / 1024

print(f"Forward pass time: {elapsed:.6f}s")
print(f"Total parameters: {params}")
print(f"Peak memory: {peak_mem:.2f} MB")
\end{verbatim}

Both generated and reference code execute with identical inputs in the sandbox, enabling objective comparison.



\subsection{Graph Routing Logic}
In the following is illustrated the conditional routing between nodes. The graph implements the following edges:

\begin{verbatim}
START → route_question
route_question → {search_relevant_sources (code), 
                  generate_query (web), 
                  generate_academic_query (academic)}
search_relevant_sources → generate
generate → collect_feedback (interrupt)
collect_feedback → process_feedback
process_feedback → {regenerate: generate,
                     approve: END,
                     evaluate: code_normalization,
                     execute: check_code_sandbox}
check_code_sandbox → reflection → collect_feedback
\end{verbatim}

Conditional edges enable dynamic routing based on state values. The feedback loop continues until the user approves or maximum iterations (default: 3) are reached.


The system implements multi-level error recovery:

\begin{enumerate}
  \item \textbf{Parser Level}: JSON parse fails → markdown extraction → raw text fallback
  \item \textbf{Agent Level}: Validate error field, code length, file markers before acceptance
  \item \textbf{Sandbox Level}: Timeout handling, missing package fallback, import error reporting
\end{enumerate}

All errors return structured dictionaries (never raw exceptions), ensuring system stability. Failed agent outputs are logged but do not terminate the workflow: remaining agents continue execution.
